\section{Object 6}
\markright{Object 6}

A helix/spiral was defined using pitch, height, and diameter parameters to generate the coil path, and a circular profile was swept along this path to form the first coil. The path and sweep were then rotated 180° and swept again to produce the second, interlocking coil.

\begin{figure}[H]
	\centering

	\begin{subfigure}[b]{0.49\textwidth}
		\centering
		\includegraphics[width=\textwidth, keepaspectratio]{lab-02/assets/object-6/1.png}
		\caption{Helix/spiral parameter dialog}
	\end{subfigure}%
	\hfill
	\begin{subfigure}[b]{0.49\textwidth}
		\centering
		\includegraphics[width=\textwidth, keepaspectratio]{lab-02/assets/object-6/2.png}
		\caption{Feature tree showing the generated helix curve}
	\end{subfigure}%
	\vspace{0.5em}
	\begin{subfigure}[b]{0.49\textwidth}
		\centering
		\includegraphics[width=\textwidth, keepaspectratio]{lab-02/assets/object-6/3.png}
		\caption{Sweeping the circular profile along the path}
	\end{subfigure}%
	\hfill
	\begin{subfigure}[b]{0.49\textwidth}
		\centering
		\includegraphics[width=\textwidth, keepaspectratio]{lab-02/assets/object-6/4.png}
		\caption{Rotating the path 180° for the second coil}
	\end{subfigure}%
	\vspace{0.5em}
	\begin{subfigure}[b]{\textwidth}
		\centering
		\includegraphics[width=\textwidth, keepaspectratio]{lab-02/assets/object-6/5.png}
		\caption{Final result}
	\end{subfigure}

	\caption{Snapshots while Designing Object 6 in Solidworks}
	\label{lab02_object_6}
\end{figure}
